\section{Methods}

\subsection{Setting and corpus}

The study examines public discussions from a specialist laboratory-automation forum covering liquid handling, integrated devices, schedulers, vendor software, open-source tools, and automation practice. The pilot contains 55 purposively selected discussions dated from 2022 through 2026. Selection sought variation across platforms, workflow layers, incidents, design questions, documentation and training requests, product reports, vendor participation, open-source development, and positive community outcomes. The corpus was designed for conceptual coverage and mechanism development. It does not support prevalence estimates.

The released analytical record excludes usernames and profile attributes. Short quotations are anonymized and linked to their public discussion in the non-anonymous research release. We did not redistribute a verbatim corpus. Public outcomes were coded as confirmed, partial, open, mixed, product-reported, or migrated. Migration records cases in which substantive work moved to direct messages, private support, Slack, meetings, or local implementation. Public silence was not interpreted as operational non-resolution.

\subsection{Units of analysis}

The thread served as the discovery unit. Fourteen deliberately difficult threads were segmented into 45 episodes because one discussion can move among interface access, configuration, physical representation, observability, recovery, validation, and public artifact creation. Each episode records the initiating problem, lifecycle stage, principal technical construct, ecosystem modifiers, evidence form, detectability, consequence, public outcome, artifact, counterexample status, and confidence.

The episode structure supports the CSCW analysis by separating forms of collaborative work that would otherwise collapse into one technical topic. A version-control discussion can contain configuration diagnosis, recovery concerns, and public bridge development. A scheduler discussion can contain capability comparison, driver-access negotiation, requirements elicitation, and community documentation. A partial-execution incident can contain preflight design, runtime evidence, material reconciliation, and recovery policy.

\subsection{Coding and analysis}

The coding model contains ten constructs. Two describe ecosystem knowledge and support. Three describe interfaces and representations. Three describe runtime coordination. One describes evaluation semantics. One describes AI context and physical feedback. Each construct has inclusion, exclusion, primary-code, evidence, adjacency, and boundary rules. Positive cases and counterexamples were retained.

All coding was performed by one analyst. The study reports no inter-rater statistic. A three-thread process-validation exercise identified ambiguities and produced a primary-code tie-break, explicit read scopes, a formal episode unit, expected episode counts, and coherence checks. A blind hard-case projection is available for a future independent pass. The answer-bearing key is stored separately and was not used as a blind instrument.

The CSCW analysis proceeded through comparison across episodes and negative cases. We examined how missing context was elicited, how participants established or contested evidence, how repair work connected software and physical state, how vendor and community roles interacted, how discussions migrated beyond the public record, and which practices converted a local answer into a durable artifact. Quantitative summaries were used to locate recurrent patterns and contrasting cases, not to infer a population distribution.

\subsection{Ethics and anonymity}

The analysis concerns public professional discussions. Data handling followed minimization, contextual quotation, and contributor de-identification. The public research release contains derived records and source links but no contributor handles or private correspondence. The review artifact for this manuscript must remove repository history, author metadata, contact details, commit authorship, and links that identify the research team. Source discussions can remain cited where their presence does not reveal authorship; the artifact containing the coded release must be anonymously mirrored for review.
